% ==========================================
% SECTION 1: TASK 1 - EXPLORATION
% ==========================================
\section{Task 1: Exploration}
\begin{frame}{Task 1: Autonomous Map Generation}

    \begin{block}{Objective}
        Perform a \textbf{full autonomous exploration} of \textbf{an unknown environment} to generate a consistent 2D occupancy grid map using SLAM.
    \end{block}

    \vspace{0.1cm}

    \begin{columns}
        \begin{column}{0.55\textwidth}
            \textbf{Environmental Challenges}
            \par
            The testing environment presents specific difficulties that standard configuration could not handle:
            \vspace{0.2cm}
            \begin{itemize}
                \item \textbf{Complex Geometry:} Narrow passages and obstacles requiring precise maneuvering.
                \item \textbf{Lighting Artifacts:} Strong shadows created false positive obstacles in the LIDAR readings (ghost objects).
            \end{itemize}
        \end{column}

        \begin{column}{0.45\textwidth}
            \begin{figure}
                \centering
                \includegraphics[width=0.6\linewidth]{kitchen_table.png}
                \caption{Environment in Gazebo 11.}
            \end{figure}
        \end{column}
    \end{columns}

\end{frame}


\begin{frame}{Task 1: Exploration Strategy}

    To achieve efficient coverage, we integrated a C++ node based on \texttt{m-explore}, enhancing the standard frontier exploration logic.

    \vspace{0.3cm}

    \begin{exampleblock}{Optimized Frontier Selection Logic}
        Instead of a greedy approach, we tuned the cost function parameters to prioritize frontiers that are:
        \begin{enumerate}
            \item \textbf{Closer} to the robot's current pose (minimizing travel time).
            \item Located in \textbf{larger unexplored voids} (maximizing information gain).
        \end{enumerate}
    \end{exampleblock}

    \vspace{0.3cm}
    
    \textbf{Completion Workflow:}
    \par
    Upon exhaustion of the whitelist frontiers, the robot automatically navigates back to the \textbf{Home Position} and triggers the map saver, storing the result in the \texttt{maps/} directory.

\end{frame}


\begin{frame}{Task 1: Safety Configuration}

    \begin{columns}
        \begin{column}{0.5\textwidth}
            \begin{alertblock}{Safety Measure: Pre-Exploration Configuration}
                Before moving, the system executes a mandatory \textbf{Tuck Arm} procedure. Folding the arm into a safe configuration minimizes the robot's footprint, preventing collisions with low obstacles or door frames during autonomous navigation.
            \end{alertblock}
        \end{column}

        \begin{column}{0.45\textwidth}
            \begin{figure}
                \centering
                \includegraphics[width=0.8\linewidth]{SafePose.png}
                \caption{Tuck Arm Configuration for Safe Navigation}
            \end{figure}
        \end{column}
    \end{columns}
    
    \vspace{0.5cm}
    This configuration is also considered in the navigation for Task 2 and Task 3 as safe Pose.

    

\end{frame}



\begin{frame}{Task 1: System Configuration}

    \begin{columns}
        \begin{column}{0.55\textwidth}
            \textbf{Sensor \& Navigation Tuning}
            \par
            We refined the navigation stack parameters to address the environmental noise:
            \begin{itemize}
                \item \textbf{LIDAR Raytracing:} Adjusted strictly to filter out shadow reflections and prevent "ghost obstacles" on the map.
                \item \textbf{Motion Profile:} Smoothed acceleration limits to ensure stable mapping and prevent odometry drift.
                \item \textbf{Cost Map Parameters:} %TODO - Add this part
            \end{itemize}
        \end{column}

        \begin{column}{0.45\textwidth}
            \begin{figure}
                \centering
                \includegraphics[width=0.8\linewidth]{Costmap.png}
                \caption{Generated Occupancy Grid}
            \end{figure}
        \end{column}
    \end{columns}

\end{frame}


\begin{frame}{Task 1: Exploration Video Result}
    \begin{figure}
        \centering
        \href{https://liveunibo.sharepoint.com/sites/idughis/_layouts/15/stream.aspx?id=\%2Fsites\%2Fidughis\%2FShared\%20Documents\%2FGeneral\%2FMagistrale\%2FAMR\%2FVideos\%2FTask1\%2Emp4&referrer=StreamWebApp\%2EWeb&referrerScenario=AddressBarCopied\%2Eview\%2Efe3d0b40\%2D24c3\%2D4d51\%2Dafda\%2D700f9ff0e2a9}{
            \includegraphics[width=0.9\linewidth]{Thumb_task1.png}
        }
        \caption{Exploration Task Execution (Click to play in browser)}
    \end{figure}
\end{frame}